%
% This work is licensed under a Creative Commons Attribution-ShareAlike 4.0 International License.
%

\documentclass[10pt, letterpaper, twoside, openany]{article}

\title{Introduction to Quantum Mechanics -- Worksheet 2}
\author{Joseph D. MacMillan}
\date{}


\usepackage{graphicx}
\usepackage{color} 
\usepackage{amsmath, amssymb}
\usepackage{libertinus}
\usepackage{microtype}
\usepackage{layout}
\usepackage[most]{tcolorbox}
\tcbuselibrary{skins,breakable}
\usepackage[hang, small, bf]{caption}
\captionsetup[table]{position=top}
\captionsetup[figure]{position=bottom}
\usepackage[hyphens]{url}
\usepackage[hidelinks]{hyperref}
\hypersetup{pdftitle={Introduction to Quantum Mechanics -- Worksheet 2}}
\hypersetup{pdfauthor={Joseph D. MacMillan}}
\urlstyle{rm}
\usepackage[shortlabels]{enumitem}
\usepackage[margin=1in]{geometry}
\usepackage{fancyhdr}




\newcommand{\uu}{\symbfup{u}}
\newcommand{\grad}{\symbfup{\nabla}}
\newcommand{\curl}{\symbfup{\nabla} \times}
\newcommand{\dfdx}[2]{\frac{\partial {#1}}{\partial {#2}}}
\newcommand{\ddfdx}[2]{\frac{\partial^2 {#1}}{\partial {#2}^2}}
\newcommand{\vort}{\symbfup{\omega}}
\renewcommand\vec{\symbfup}
\newcommand{\unit}[1]{\hat{\vec{#1}}}
\newcommand{\U}{\mathbb{U}}
\newcommand{\V}{\mathbb{V}}
\newcommand{\LL}{\mathbb{L}^2}
\newcommand{\LLLL}{\mathbb{L}^4}

\newcommand{\ket}[1]{| #1 \rangle}
\newcommand{\bra}[1]{\langle #1 |}
\newcommand{\braket}[2]{\langle #1 | #2 \rangle}


%\DeclareMathAlphabet{\mathbfsf}{\encodingdefault}{\sfdefault}{bx}{sl}
\newcommand{\tens}[1]{\mathbfsfit{#1}}


\newcounter{example}[section]
\def\theexample{\thechapter.\arabic{example}}
\newenvironment{example}[1][ ]{\refstepcounter{example}
\begin{tcolorbox}[breakable, sharp corners, boxrule = 0pt, frame empty, opacityframe=0, parbox=false]
\textbf{Example \theexample \ -- #1.}
}
{ 
\end{tcolorbox}
}

\newenvironment{theorem}[1][ ]{
\begin{tcolorbox}[breakable, sharp corners, boxrule = 0pt, frame empty, opacityframe=0, parbox=false]
\textbf{#1.}
}
{ 
\end{tcolorbox}
}

\newcounter{problem}[section]
\def\theproblem{\thechapter.\arabic{problem}}
\newenvironment{problem}[1][ ]{\refstepcounter{problem} \noindent \textbf{Problem \theproblem \ -- #1.}}{\vspace{0.1in}}

\renewcommand{\arraystretch}{2.5}

\begin{document}

\pagestyle{fancy}
\fancyhead[L]{\emph{Introduction to Quantum Mechanics}}
\fancyhead[R]{Worksheet 2} 

\Large \noindent Worksheet 2

\noindent \LARGE \textbf{Stern-Gerlach and Probabilities}

\normalsize

\vspace{0.2in}

\noindent Exploring the Stern-Gerlach Experiment  \bullet \  Normalization   \bullet \  Probabilities  

\vspace{0.1in}

\noindent \rule{\textwidth}{0.5pt}

\vspace{0.2in}

\begin{enumerate}
\item This question uses the SPINS simulation which you can find online at \url{https://physics.weber.edu/schroeder/software/Spins.html}.  Set up the experiment so that you have an oven, a Stern-Gerlach $S_z$ device, and two counters, one for each output. Click on the number $1$ on the oven; this will set each particle coming out to be in the initial state
\[
\ket{\psi_1} = \frac{1}{\sqrt{2}} \ket{+} + \frac{i}{\sqrt{2}} \ket{-}.
\]
\begin{enumerate}
\item First, predict what the experiment will measure:  what are the probabilities $\mathcal{P}_+$ and $\mathcal{P}_-$ for measuring the particle to be spin up or spin down, respectively?  Show your calculation.

\vspace{2in}

\item Okay, now go ahead and click the 10k button to send in 10,000 particles.  Are your predictions correct?  Are they \emph{perfectly} correct?

\vspace{1in}

\item Next, replace the $S_z$ device with one oriented along the $x$ direction.  Once again, make a prediction: what are the probabilities for measuring $\mathcal{P}_{+x}$ and $\mathcal{P}_{-x}$?  Test your prediction once you've done the calculation.

\newpage

\item Finally, replace the $S_x$ device with one oriented along the $y$ direction.  What are the probabilities for measuring $\mathcal{P}_{+y}$ and $\mathcal{P}_{-y}$?  Test your predictions.  Any observations to make here?

\end{enumerate}

\vspace{3in}

\item Now for something a little trickier.  One the SPINS experiment, keep the single Stern-Gerlach device (any orientation is fine), but click on the number $3$.  This one sends out particles all in the same state $\ket{\psi_3}$, but we don't know exactly what that state is.  Our job is to figure it out.
\begin{enumerate}
\item First, run your experiments: using all three orientations for the SG device, figure out the probabilities for spin up and spin down for each and record them below.

\begin{center}
\begin{tabular}{ p{1in} | p{1in} | p{1in} | p{1in}}
Probabilities & $S_x$ & $S_y$ & $S_z$ \\
\hline
Spin up & & & \\
\hline
Spin down & & & \\
\hline
\end{tabular}
\end{center}


\item Write the unknown state $\ket{\psi_3}$ as 
\[
\ket{\psi_3} = a \ket{+} + b\ket{-},
\]
where $a$ and $b$ are unknown \emph{complex} constants.  This state is written in the $S_z$ basis, so start with your results for that orientation in table.  Calculate $\mathcal{P}_+$ and $\mathcal{P}_-$ and use those results to find out what $|a|$ and $|b|$ are.  

\newpage

\item Unfortunately, we only know the magnitudes of those complex numbers $a$ and $b$.  Since we know (or will soon) that an overall phase has no physical meaning, we can write the state now as 
\[
\ket{\psi_3} = |a| \ket{+} + |b| e^{-\phi}\ket{-},
\]
where $\phi$ is the phase -- our last unknown.  Calculate  $\mathcal{P}_{+x}$ and $\mathcal{P}_{-x}$ and use your results from the $S_x$ experiment to find what $\phi$ is (careful -- there are multiple possibilities it could be!).  \emph{Hint:} You might need to use the relationship
\[
\cos\phi = \frac{e^{i\phi} + e^{-i\phi}}{2}.
\]

\vspace{3in}

\item Okay, one last problem -- we don't know the sign of the phase $\phi$ (it could be positive or negative).  To resolve this last ambiguity, we need to use our results for $S_y$ -- go ahead, and write the final ket of the state $\ket{\psi_3}$. 

\end{enumerate}
\end{enumerate}

\end{document}
